\chapter{Caucher Birkar}\label{ch2018caucherbirkar:caucher-birkar-fields-medal-2018-a-comprehensive-report-on-his-work-and-impact}

\begin{center}\rule{0.5\linewidth}{0.5pt}\end{center}

\textbf{Prepared:} July 2026

\textbf{Subject:} Caucher Birkar (b.\ July 1978, Marivan, Kurdistan Province, Iran)

\textbf{Position:} Professor, Yau Mathematical Sciences Center (YMSC) and BIMSA, Tsinghua University; Professor, University of Cambridge

\textbf{Doctoral Advisor:} Ivan Fesenko (University of Nottingham)

\textbf{Fields Medal Citation:} \emph{``For the proof of the boundedness of Fano varieties and for contributions to the minimal model program.''}

\begin{center}\rule{0.5\linewidth}{0.5pt}\end{center}

\begin{prereqs}
Algebraic geometry (projective varieties, divisors, linear systems, birational morphisms), basic commutative algebra (local rings, completion, normalization).
\end{prereqs}

\begin{objectives}
Explain the boundedness of Fano varieties (the BAB conjecture); describe the minimal model program in higher dimensions and the existence of minimal models; understand the construction of moduli spaces of Fano varieties.
\end{objectives}

\section{Problem History}

The classification of algebraic varieties up to birational equivalence is a central goal of algebraic geometry. In dimension 1 (curves), classification was completed by Riemann in the 19th century using the genus as a complete invariant. In dimension 2 (surfaces), the Italian school (Enriques, Castelnuovo, Kodaira) classified surfaces into four classes: general type, ruled, Abelian, and rational. For dimension 3 and higher, the \textbf{Minimal Model Program} (MMP), initiated by Shigefumi Mori in the 1980s, provides the framework for classification.

The MMP proceeds by elementary operations---contractions, flips, and divisorial extractions---to find a ``minimal'' representative in each birational class. The two possible outputs are a \textbf{minimal model} (with nef canonical divisor $K_X$) or a \textbf{Mori fiber space} (a Fano fibration). A fundamental question is whether the \textbf{canonical ring} $R(X, K_X) = \bigoplus_{m \ge 0} H^0(X, \mathcal{O}_X(mK_X))$ is finitely generated. If so, the variety has a well-defined birational model that can be studied systematically.

The Borisov--Alexeev--Borisov (BAB) conjecture, formulated in the 1990s, predicts that for fixed dimension $n$ and a fixed DCC set of coefficients $D \subset [0,1]$, the set of all $n$-dimensional Fano varieties with $\epsilon$-log terminal singularities and coefficients in $D$ is \textbf{bounded}---parametrized by a finite number of algebraic families. Before Birkar, it was not even known whether the number of deformation types of terminal Fano varieties was finite. The conjecture is a vast generalization of the boundedness of curves of genus $g$ and the classification of surfaces.

The \textbf{Shokurov boundedness of complements conjecture} predicted that for fixed $n$ and DCC $D$, the set of all complements on $n$-dimensional Fano varieties with coefficients in $D$ is bounded. A complement is a divisor $D' \sim -mK_X$ such that $(X, D')$ is log canonical---this constrains the possible singularities and provides a powerful tool for studying Fano varieties.

The \textbf{conjecture of Borisov} on the existence of complements predicted that for every Fano variety of log index at least $n$, there exists a complement. This was a key step toward proving the BAB conjecture and understanding the structure of Fano varieties. The boundedness of complements and the existence of complements together provide a complete picture of the structure of Fano varieties.

The \textbf{Koll\'ar--Shepherd-Barron--Alexeev theorem} (1988) provides the foundation for moduli theory of pairs: the moduli space of stable pairs $(X, \Delta)$ with $\Delta$ in a fixed numerical class is a proper algebraic stack. The boundedness of Fano varieties would imply the existence of moduli spaces of K-stable Fano varieties, connecting birational geometry to differential geometry and mirror symmetry.

The \textbf{K\"ahler--Einstein metric problem} asks whether a Fano variety admits a K\"ahler--Einstein metric. The Yau--Tian--Donaldson conjecture (proved by Chen--Donaldson--Sun and Tian in 2015) states that this is equivalent to K-stability of the variety. The boundedness of Fano varieties implies that there are only finitely many deformation types of K-stable Fano varieties in each dimension, which is a powerful constraint on the existence of K\"ahler--Einstein metrics.

\section{How Was It Solved}

Birkar's approach \cite{birkar2018} built on a series of foundational results in birational geometry, culminating in the proof of the BAB conjecture. The key innovation was the development of \textbf{generalized pairs}, a new framework that extends the classical theory of pairs to situations where the ``boundary'' includes a virtual nef divisor from a higher-dimensional space.

\textbf{The BCHM Theorem (2010).} Birkar, Cascini, Hacon, and McKernan proved that the MMP works for varieties of log general type: minimal models exist, canonical rings are finitely generated, and flips exist for klt pairs in all dimensions. This established the foundational results the MMP requires and resolved a conjecture of Zariski from the 1940s. The proof uses the base point free theorem and the contraction theorem to show that the MMP terminates, with the key technical innovation being the use of b-semiampleness and the theory of movable divisors. The termination of the MMP is one of the deepest results in algebraic geometry, and its proof requires a delicate balance of algebraic, analytic, and geometric techniques.

The BCHM theorem builds on the work of many mathematicians: Mori's bend-and-break technique for constructing rational curves on Fano varieties, the cone theorem and contraction theorem from the 1980s, and the theory of minimal models developed by Kawamata, Viehweg, and others in the 1990s. The key innovation of BCHM was to develop new techniques for controlling the termination of the MMP, particularly in situations where the standard arguments fail. The proof uses the theory of log canonical thresholds, the minimal model program with scaling, and a careful analysis of the behavior of the canonical divisor under birational transformations.

The existence of flips in all dimensions is particularly remarkable: a flip is a birational modification that replaces a negative extremal ray with a positive one, and constructing such modifications requires solving a nonlinear PDE (the K\"ahler--Einstein equation) on the flipped variety. The BCHM theorem provides an algebraic construction of flips that avoids the analytic difficulties, using only techniques from commutative algebra and algebraic geometry.

The \textbf{flips} that appear in the MMP are obtained by contracting an extremal ray of the cone of curves and then ``flipping'' the contracted variety to obtain a new variety with better properties. The existence of flips was proved by Mori (1988) in dimension 3 and by Birkar--Cascini--Hacon--McKernan in all dimensions. The flip is a birational modification that replaces a negative extremal ray with a positive one, analogous to the flop in symplectic geometry.

\textbf{Anti-pluricanonical systems (2016--2019).} Birkar proved that for fixed dimension $n$, there exists $m = m(n)$ such that the anti-pluricanonical system $|-mK_X|$ is non-empty and birational for every $n$-dimensional terminal Fano variety. The proof uses volume bounds ($\operatorname{vol}(-K_X) \le C(n)$) and multiplicity estimates, combined with inductive MMP arguments. The volume bound is the key estimate: it shows that Fano varieties cannot have arbitrarily large anticanonical volume, which constrains their geometry and prevents them from ``escaping to infinity'' in the moduli space.

The volume bound is proved using a delicate induction on dimension. The key step is to show that the anticanonical volume of an $n$-dimensional Fano variety is bounded in terms of the anticanonical volumes of its lower-dimensional fibers. This requires a careful analysis of the fibration structure of Fano varieties and the behavior of the canonical divisor under specialization.

\textbf{Boundedness of complements.} Birkar solved Shokurov's conjecture: for fixed $n$ and DCC $D$, the set of all complements on $n$-dimensional Fano varieties with coefficients in $D$ is bounded. The proof uses a careful analysis of the singularities of Fano varieties and the theory of generalized pairs. The boundedness of complements is the crucial input for the proof of the BAB conjecture.

A \textbf{complement} on a Fano variety $X$ is a divisor $D' \sim -mK_X$ such that $(X, D')$ is log canonical and $D'$ has coefficients in a fixed DCC set $D$. The boundedness of complements means that there are only finitely many deformation types of complements, which constrains the possible singularities and geometry of Fano varieties. The existence of complements was conjectured by Borisov and proved by Birkar: for every Fano variety of log index at least $n$, there exists a complement. Together, the boundedness and existence of complements provide a complete picture of the structure of Fano varieties.

\textbf{The proof of the BAB conjecture (2021).} Combining anti-pluricanonical systems, boundedness of complements, and the theory of generalized pairs, Birkar proved the BAB conjecture. The proof, published in the \emph{Annals of Mathematics}, runs over 100 pages and reduces the problem through birational boundedness to a finite number of cases. The final step uses a limiting argument to show that the boundedness holds for all Fano varieties, not just those with bounded anticanonical volume.

\section{Impact of This Discovery}

The BCHM theorem is now a foundational tool in algebraic geometry. Finite generation of canonical rings, proved for all dimensions, resolved a conjecture of Zariski from the 1940s. The boundedness of Fano varieties enabled the construction of moduli spaces of K-stable Fano varieties by Chenyang Xu and others, connecting birational geometry to differential geometry (K\"ahler--Einstein metrics) and mirror symmetry.

Birkar's results imply the Jordan property for Cremona groups: for fixed $n$, there is a constant $J(n)$ such that every finite subgroup of $\operatorname{Bir}(\mathbb{P}^n)$ has an abelian normal subgroup of index at most $J(n)$. This provides a uniform bound on the structure of finite groups of birational automorphisms.

The theory of generalized pairs, developed extensively by Birkar, is now a standard technique used worldwide in birational geometry, enabling the study of singularities and the construction of moduli spaces. The construction of moduli spaces of K-stable Fano varieties has applications to differential geometry, mirror symmetry, and the geometry of Calabi--Yau manifolds.

The \textbf{BAB conjecture} (Borisov--Alexeev--Borisov) predicts that for fixed dimension $n$ and a fixed DCC set of coefficients $D \subset [0,1]$, the set of all $n$-dimensional Fano varieties with $\epsilon$-log terminal singularities and coefficients in $D$ is bounded. Birkar's proof establishes this by showing that the anticanonical volume $\operatorname{vol}(-K_X)$ is bounded in terms of $n$ and $D$, which prevents the Fano varieties from ``escaping to infinity'' in the moduli space. The boundedness then follows from the theory of Hilbert schemes and the existence of ample divisors on bounded families.

The \textbf{minimal model program} (MMP) is a systematic procedure for constructing minimal models of algebraic varieties. The MMP proceeds by a sequence of birational transformations---contractions, flips, and divisorial extractions---that simplify the variety while preserving its birational class. The two possible outcomes are a minimal model (with nef canonical divisor) or a Mori fiber space (a Fano fibration). The BCHM theorem establishes that the MMP works for varieties of log general type, which means that minimal models exist and the canonical ring is finitely generated.

The theory of \textbf{generalized pairs} extends the classical theory of pairs $(X, \Delta)$ to situations where the ``boundary'' $\Delta$ includes a virtual nef divisor from a higher-dimensional space. A generalized pair $(X, \Delta + M)$ consists of a variety $X$, a boundary divisor $\Delta$, and a nef $\mathbb{Q}$-divisor $M$ on a variety $Y$ that dominates $X$. The theory of generalized pairs has become essential for the MMP, as many constructions in the MMP require passing to a higher-dimensional space to perform birational transformations.

\section{Modern Views and Improvements}

Birkar continues to develop the theory of generalized pairs, proving boundedness and volume results for log Calabi--Yau fibrations. His recent work on positive characteristic geometry completes the MMP for 3-folds in large characteristic, extending the MMP to a new setting.

Connections to physics have emerged: Birkar and Lee (2025) apply birational geometry to supergravity, using boundedness to derive bounds on the particle spectrum of 6-dimensional theories. This represents a novel application of pure mathematics to theoretical physics.

The construction of moduli spaces of K-stable Fano varieties, building on the boundedness result, is an active area connecting algebraic geometry, differential geometry, and symplectic geometry. The K-stability condition, introduced by Tian and Donaldson, provides a differential-geometric characterization of Fano varieties that admit K\"ahler--Einstein metrics.

Current research includes the extension of Birkar's boundedness results to Fano varieties over non-algebraically closed fields, the development of moduli theory for log Fano pairs, and the study of the geometry of the moduli space itself. The interplay between Birkar's algebraic methods and the differential-geometric approach to K-stability continues to produce new insights.

The construction of moduli spaces of K-stable Fano varieties is an active area that connects Birkar's algebraic results to the differential-geometry of K\"ahler--Einstein metrics. The moduli space is expected to be a proper algebraic stack, and its geometry encodes information about the deformation theory of Fano varieties. Recent work by Xu, Chi, and others has established the existence of the moduli space in several cases, but the general theory remains under development.

Another direction is the extension of Birkar's results to positive characteristic. The MMP in positive characteristic is much more difficult than in characteristic zero, because the Frobenius morphism introduces new complications. Birkar's recent work on positive characteristic geometry has established the MMP for 3-folds in large characteristic, but the general case remains open. The development of new techniques for handling the Frobenius morphism is a key challenge in this area.

The connection between Birkar's work and mirror symmetry is another active direction. The mirror symmetry correspondence predicts a duality between Fano varieties and Landau--Ginzburg models, and the boundedness of Fano varieties implies constraints on the structure of the mirror. Understanding this connection could lead to new results on both sides of the mirror symmetry correspondence.

The study of Fano varieties in positive characteristic presents new challenges. The Frobenius morphism introduces complications that are absent in characteristic zero, and the MMP in positive characteristic requires fundamentally different techniques. Birkar's recent work in collaboration with others has established the MMP for 3-folds in large characteristic, but the general case remains open.

Another direction is the study of the moduli space of K-stable Fano varieties. The existence of this moduli space, building on Birkar's boundedness result, connects birational geometry to differential geometry and mirror symmetry. The geometry of the moduli space itself---its singularities, its compactification, and its relationship to the moduli space of canonically polarized varieties---is an active area of research.

The development of the theory of K-stability has been closely connected to Birkar's work. K-stability, introduced by Tian and Donaldson, provides a differential-geometric characterization of Fano varieties that admit K\"ahler--Einstein metrics. The boundedness of Fano varieties implies that there are only finitely many deformation types of K-stable Fano varieties in each dimension, which is a powerful constraint on the existence of K\"ahler--Einstein metrics. The construction of moduli spaces of K-stable Fano varieties is an active area that connects algebraic geometry, differential geometry, and symplectic geometry.

The application of Birkar's results to the study of birational automorphism groups has led to new results on the structure of Cremona groups. The Jordan property for Cremona groups, which follows from the boundedness of Fano varieties, provides a uniform bound on the structure of finite subgroups of birational automorphisms. This result has implications for the study of rationality questions and for the classification of algebraic varieties up to birational equivalence.

The development of the theory of complements, introduced by Shokurov and developed extensively by Birkar, has been a powerful tool for studying Fano varieties. A complement on a Fano variety $X$ is a divisor $D' \sim -mK_X$ such that $(X, D')$ is log canonical and $D'$ has coefficients in a fixed DCC set. The boundedness of complements provides a uniform bound on the possible singularities and geometry of Fano varieties, which is the key input for the proof of the BAB conjecture. The theory of complements continues to be developed and has applications to other problems in birational geometry, including the study of log Calabi-Yau fibrations and the classification of Fano varieties of small dimension.

The study of log Calabi-Yau pairs in birational geometry has been influenced by Birkar's work. A log Calabi-Yau pair $(X, \Delta)$ consists of a variety $X$ and a boundary divisor $\Delta$ such that $K_X + \Delta \sim 0$. The structure theory of log Calabi-Yau pairs, including the abundance theorem and the classification of log Calabi-Yau fibrations, has been developed using Birkar's techniques. The connection between log Calabi-Yau pairs and mirror symmetry, through the theory of SYZ fibrations, is an active area that builds on Birkar's foundational work in birational geometry.

The application of Birkar's results to the study of algebraic cycles on Fano varieties has led to new results on the Hodge conjecture for special classes of Fano varieties. The Hodge conjecture predicts that every Hodge class on a smooth projective variety is a rational linear combination of algebraic cycles, and the boundedness of Fano varieties provides constraints on the possible Hodge classes. The development of new techniques for studying Hodge classes on Fano varieties is an active area that builds on Birkar's boundedness results.

\begin{center}\rule{0.5\linewidth}{0.5pt}\end{center}

\section*{Summary}
\addcontentsline{toc}{section}{Summary}

Birkar proved the boundedness of Fano varieties (the BAB conjecture) and established foundational results for the minimal model program, including finite generation of canonical rings, existence of minimal models for log general type, and the development of generalized pairs.
